\documentclass[11pt]{article}

\usepackage[margin=1in]{geometry}
\usepackage[T1]{fontenc}
\usepackage{lmodern}
\usepackage{booktabs}
\usepackage{tabularx}
\usepackage{array}
\usepackage{enumitem}
\usepackage{xcolor}
\usepackage{hyperref}
\usepackage{titlesec}

\definecolor{accent}{RGB}{30,60,110}
\hypersetup{colorlinks=true, urlcolor=accent, linkcolor=accent}
\titleformat{\section}{\large\bfseries\color{accent}}{\thesection}{0.6em}{}
\titleformat{\subsection}{\normalsize\bfseries}{\thesubsection}{0.6em}{}

\newcolumntype{L}{>{\raggedright\arraybackslash}X}
\newcommand{\checkbox}{\raisebox{-0.5pt}{\framebox(7,7){}}}
\newlist{todo}{itemize}{2}
\setlist[todo]{label=\checkbox, leftmargin=1.6em, itemsep=2pt}

\setlength{\parindent}{0pt}
\setlength{\parskip}{6pt}

\begin{document}

\begin{center}
  {\LARGE\bfseries REIT FFO Comparability Pipeline}\\[4pt]
  {\large Project Plan}\\[8pt]
  Mohammad Natsheh \quad\textbar\quad MRED, Clemson University \quad\textbar\quad October 2026
\end{center}

\section{Project Summary}

Funds from operations (FFO) is the main earnings measure for REITs, but the
company-defined versions (Core FFO, Normalized FFO, CFFO, FFO as Adjusted,
AFFO) are not comparable across peers. Each REIT adds back different items,
companies change their own definitions over time, and the reconciliations sit
in quarterly filings with different layouts.

\textbf{One-sentence description.} Our application helps REIT equity analysts
and real estate investors make apples-to-apples valuation and peer-comparison
decisions across all U.S.\ publicly traded multifamily REITs by using an AI
agent that monitors SEC EDGAR daily to track which REITs are listed, collect
every quarterly earnings filing since 2018, and extract each FFO
reconciliation, to produce a continuously updated, restated FFO history on one
common definition, with source-page citations and plain-language explanations
of how each REIT's reported metrics differ.

The pipeline is written in Python on the Claude API and:
\begin{enumerate}[itemsep=1pt]
  \item \textbf{tracks the universe}: finds every listed multifamily REIT,
        including those listed at any point since 2018 that have since been
        acquired, merged or liquidated;
  \item \textbf{collects} each quarterly earnings document from SEC EDGAR;
  \item \textbf{extracts} every line of each FFO reconciliation into structured data;
  \item \textbf{restates} every company-quarter to one fixed common definition,
        applied retroactively (the arithmetic runs in code, not in the AI);
  \item \textbf{explains} how each REIT's reported metrics differ, with page citations; and
  \item \textbf{updates itself daily}, picking up new quarters, new REITs and exits.
\end{enumerate}

\begin{center}
\small
\fbox{\parbox{0.92\textwidth}{\centering
\textbf{Track universe} $\rightarrow$ \textbf{Collect} (EDGAR API) $\rightarrow$
\textbf{Extract} (AI) $\rightarrow$ \textbf{Restate} (Python) $\rightarrow$
\textbf{Explain and report} \\[2pt] \textit{runs daily on a schedule}}}
\end{center}

\section{Scope}

\begin{itemize}[itemsep=2pt]
  \item \textbf{Companies:} every U.S.\ exchange-listed REIT that meets the
        multifamily rule (Phase 4), at any point from 2018 to today.
  \item \textbf{Period:} Q1 2018 to the latest reported quarter, then every new
        quarter as it is filed. Roughly 35 quarters of history.
  \item \textbf{Metrics:} NAREIT FFO, each company's own ``core'' metric, and AFFO.
  \item \textbf{Common definition:} one fixed definition, today's, applied to
        every quarter regardless of what definition the company used at the time.
  \item \textbf{Disclosure limits:} where an older filing combines or omits an
        item the definition needs, the pipeline searches the 10-Q notes; if the
        item still cannot be separated, the quarter is flagged ``restated with
        limits'' rather than guessed.
\end{itemize}

\subsection*{Current universe (October 2026)}

{\small\noindent
\begin{tabularx}{\textwidth}{@{}l >{\raggedright\arraybackslash}p{3.9cm} >{\raggedright\arraybackslash}p{3.0cm} L@{}}
\toprule
\textbf{Ticker} & \textbf{Company} & \textbf{Markets} & \textbf{``Core'' metric / note} \\
\midrule
VMRK & Vivmark Residential & Coastal & Formed Aug.\ 2026 from AvalonBay (Core FFO) and Equity Residential (Normalized FFO); history kept as two companies before the merger \\
ESS  & Essex Property Trust & West Coast & Core FFO \\
UDR  & UDR & Mixed & FFO as Adjusted \\
MAA  & Mid-America Apartment Communities & Sunbelt & Core FFO, Core AFFO \\
CPT  & Camden Property Trust & Sunbelt & Core FFO, Core AFFO \\
IRT  & Independence Realty Trust & Sunbelt/Midwest & CFFO; acquiring Centerspace \\
NXRT & NexPoint Residential Trust & Sunbelt & Core FFO \\
BRT  & BRT Apartments & Sunbelt & No core metric (FFO, AFFO) \\
CSR  & Centerspace & Midwest/Mountain & Core FFO; being acquired by IRT \\
CLPR & Clipper Realty & New York City & AFFO; inclusion decided by the multifamily rule \\
\bottomrule
\end{tabularx}}

\subsection*{Former REITs in the 2018+ history (examples)}
Veris Residential (taken private 2026), Elme Communities (liquidated 2026),
Aimco (liquidating), AIR Communities (taken private 2024), Preferred Apartment
Communities and Bluerock Residential (acquired 2022). Including them avoids
\textbf{survivorship bias}: a history built only from today's survivors would
skew every comparison.

\section{Testing Strategy}

The hand-built set is the answer key the pipeline is graded against:
\begin{itemize}[itemsep=1pt]
  \item 15 supplementals collected by hand: MAA, CPT, IRT, NXRT and BRT for
        Q4 2025--Q2 2026.
  \item MAA reconciled by hand for all three quarters (Excel workbook).
  \item Spot checks: one quarter each for three or four other REITs, plus at
        least two pre-2020 quarters to test older layouts.
\end{itemize}

\section{Project Steps}

\subsection{Phase 1 --- Scope and Background}
\begin{todo}
  \item Confirm scope: all U.S.-listed multifamily REITs, Q1 2018 to present. \textit{Done.}
  \item Read the Nareit FFO White Paper and list the NAREIT-standard adjustments.
\end{todo}

\subsection{Phase 2 --- Answer Key Documents}
\begin{todo}
  \item Collect 15 supplementals by hand. \textit{Done.}
  \item Rename files consistently and store in \texttt{data/raw/}. \textit{Done.}
  \item Replace the BRT Q4 2025 file with a text-based copy from EDGAR.
  \item Record the EDGAR filing each file came from (form, date, exhibit).
\end{todo}

\subsection{Phase 3 --- Manual Baseline (MAA)}
\begin{todo}
  \item Q2 2026 reconciliation. \textit{Done (worked example).}
  \item Q4 2025 and Q1 2026 reconciliations; all checks show OK.
  \item Fill in the Adjustment Log for all three quarters.
\end{todo}

\subsection{Phase 4 --- Definitions}
\begin{todo}
  \item Build the adjustment map across the five answer-key REITs.
  \item Write the common FFO definition: which items are included and why.
  \item Write the \textbf{multifamily rule} (e.g.\ more than half of NOI or
        real estate from apartments) used to decide which REITs are in the universe.
  \item Define the ``restated with limits'' flag and when it applies.
\end{todo}

\subsection{Phase 5 --- Collection Agent (fixed list, current quarters)}
\begin{todo}
  \item Set up Python, the Claude API key and the repository. \textit{Repository done.}
  \item Write EDGAR tools: list a company's filings, list a filing's exhibits, download a file.
        Send the User-Agent header the SEC requires.
  \item Agent finds each earnings 8-K and picks the exhibit holding the FFO reconciliation.
  \item Checks: file mentions ``Funds from Operations'' and the right period end.
  \item Run on the five answer-key REITs; compare against the 15 hand-collected files.
\end{todo}

\subsection{Phase 6 --- Extraction and Restatement}
\begin{todo}
  \item Define the JSON schema for one reconciliation (lines, amounts, per-share, pages).
  \item Prompt Claude to fill it and to flag unclear items instead of guessing.
  \item Check in code that extracted lines sum to reported totals.
  \item Map line items to common categories; compute restated FFO in Python.
  \item Generate a plain-language note per company-quarter.
  \item Test on MAA against the manual baseline until it matches.
\end{todo}

\subsection{Phase 7 --- Accuracy Testing}
\begin{todo}
  \item Collection: correct filing and exhibit for every answer-key company-quarter.
  \item Extraction: match to \$0.01 per share.
  \item Error log by type: wrong document, wrong number, misclassified adjustment,
        wrong page, missed item.
\end{todo}

\subsection{Phase 8 --- Live Universe and Daily Updates}
\begin{todo}
  \item Pull all REITs (SIC code 6798) with tickers from the SEC company list.
  \item Claude classifies each candidate against the multifamily rule, citing the 10-K.
  \item New entries are proposed for review before they join the universe.
  \item Detect exits from Form 25 (delisting) and Form 15 (deregistration); keep history.
  \item Schedule a daily run with GitHub Actions: new quarters, new REITs, exits.
\end{todo}

\subsection{Phase 9 --- Backfill to Q1 2018}
\begin{todo}
  \item Identify every multifamily REIT listed at any point since 2018, including former ones.
  \item Collect and extract all quarters from Q1 2018 onward.
  \item Search 10-Q notes where reconciliations lack detail; flag what cannot be separated.
  \item Spot-check pre-2020 quarters by hand.
  \item Report the share of quarters fully restated vs.\ restated with limits.
\end{todo}

\subsection{Phase 10 --- Deliverables}
\begin{todo}
  \item Full proposal after fall break (scope, method, universe, test plan).
  \item Code repository on GitHub with a README.
  \item Restated FFO history (Excel and data files).
  \item Optional results website (e.g.\ Vercel) showing the latest comparison.
  \item Write-up of findings and accuracy results.
  \item Short resume summary of the project.
\end{todo}

\section{Data Sources}
\begin{itemize}[itemsep=1pt]
  \item SEC EDGAR via the free API (\texttt{data.sec.gov}): 8-K earnings exhibits,
        10-Q and 10-K filings, Forms 25 and 15, and the company list with SIC codes.
  \item REIT investor relations sites: quarterly supplementals (answer key).
  \item Nareit FFO White Paper (baseline definition).
  \item Claude API: classification, extraction and explanations.
\end{itemize}

\section*{References}
\small
\begin{itemize}[itemsep=1pt, label={}]
  \item \href{https://www.regardsofwallstreet.com/news/avalonbay-avb-merger-closes-vmrk-august-17}{AvalonBay--Equity Residential merger closes as Vivmark}
  \item \href{https://www.insidearbitrage.com/?p=174660}{IRT to acquire Centerspace}
  \item \href{https://www.multihousingnews.com/affinius-capital-jv-closes-3-5b-acquisition-of-veris-residential/}{Veris Residential taken private}
  \item \href{https://finviz.com/news/407432/elme-communities-announces-suspended-nyse-trading-anticipates-trading-will-begin-otc-effective-october-26-2026}{Elme Communities NYSE trading suspended}
  \item \href{https://digital-release.kxan.com/business/press-releases/cision/20260330LA22745/aimco-provides-update-on-liquidation-activities}{Aimco liquidation update}
  \item \href{https://www.sec.gov/Archives/edgar/data/0000912595/000119312526323756/maa-ex99_1.htm}{MAA Q2 2026 8-K}
\end{itemize}

\end{document}
